\chapter{Planos y Croquis de Infraestructura}

El estudio técnico contratado al Ing.\ Remigio Ojeda en 2025 consolidó por primera vez un paquete integral de planos, croquis y perfiles del sistema de agua potable. La documentación completa se conserva en formato físico y digital (\emph{anexos/B\_planos\_completos/}); a continuación se resumen los elementos indispensables para la entrega.

\section{Documentos entregados}

La revisión documental confirmó la recepción de planos generales y específicos que sirven de base para decisiones técnicas y administrativas.

\subsection*{Resumen ejecutivo}
\begin{itemize}[noitemsep, topsep=0pt]
  \item Se validó la existencia de planos firmados y numerados con la descripción de la red de distribución y la planta de tratamiento.
  \item Se dejó constancia de la ubicación del respaldo digital en Google Drive y de la copia física en carpeta rotulada ``Planos 2025''.
\end{itemize}
\textit{Nota: Los archivos en formato PDF y DWG se listan íntegramente en el anexo mencionado.}

\subsection*{Ejemplos de planos disponibles}
\begin{table}[H]
\centering
\caption{Planos representativos (ver anexo para colección completa)}
\label{tab:planos_representativos}
\begin{tabular}{p{4.5cm} p{5.2cm} p{4cm}}
\toprule
\textbf{Nombre del plano} & \textbf{Cobertura} & \textbf{Observaciones} \\
\midrule
Plano general de red & Conducción y distribución parroquial & Incluye coordenadas en formato UTM, verificación en campo 2025 \\
Plano de planta de tratamiento & Captación, filtración y caseta de cloración & Señala mejoras sugeridas para cloración y filtración lenta \\
\bottomrule
\end{tabular}
\end{table}
\textit{Los planos de detalle (válvulas, estructuras especiales y perfiles) permanecen en el anexo digital.}

\section{Croquis operativos}

Para operaciones diarias se mantienen croquis simplificados con referencias de válvulas, cajas y límites de sectores. Estos croquis sirven para intervención rápida durante emergencias.

\subsection*{Sectores priorizados}
\begin{itemize}[noitemsep, topsep=0pt]
  \item Sector Centro: válvula de control ubicada en [Coordenada/Referencia], permite aislar el casco parroquial.
  \item Sector La Campanera: croquis con puntos de riesgo por deslizamientos identificados en 2025.
\end{itemize}
\textit{Los demás sectores y subsectores se describen en los croquis completos archivados en anexos.}

\subsection*{Válvulas representativas}
\begin{table}[H]
\centering
\caption{Válvulas críticas (consultar anexos para listado total)}
\label{tab:valvulas_representativas}
\begin{tabular}{p{2.2cm} p{6.2cm} p{5.3cm}}
\toprule
\textbf{Código} & \textbf{Ubicación referencial} & \textbf{Función principal} \\
\midrule
V-01 & Entrada a planta, junto a cámara de mezcla rápida & Permite cortar alimentación general para mantenimiento de planta \\
V-03 & Intersección Tres Esquinas, Barrio Progreso & Sectoriza el barrio y reduce pérdidas durante reparaciones \\
\bottomrule
\end{tabular}
\end{table}
\textit{La numeración completa y las coordenadas precisas se encuentran en el anexo digital y en la libreta del operador.}

\section{Recomendaciones de uso}

\begin{itemize}[noitemsep, topsep=0pt]
  \item Programar jornadas de socialización interna con operador y nueva directiva para repasar planos y rutas de cierre de válvulas.
  \item Actualizar los documentos cada vez que se ejecute una ampliación, reemplazo o nueva intervención estructural.
\end{itemize}
\textit{Cualquier actualización deberá conservarse en el mismo repositorio digital y comunicarse mediante acta a la asamblea.}
